\chapter{Experiments and Discussion}
\label{cha:experiments}

\section{Experimental Setup}
\label{sec:setup}

\textbf{Hardware Platform.} Experiments were conducted on the Processing System (PS) of
the Kria KV260, an embedded UMA platform featuring a quad-core CPU running at 1.3\,GHz
and an FPGA fabric operating at 250\,MHz. The system is equipped with 4\,GB of shared
memory. Power measurements were collected using onboard sensors sampling at runtime via
the AMD \texttt{xmutil} toolchain.

\textbf{Software Stack.} The inference engine is based on a modified version of
\texttt{llama.cpp}, extended to support hardware offloading of the MatMul + SwiGLU
block. Models are executed using the Q4\_K quantization format. The FPGA accelerator was
synthesized using the Vitis high-level synthesis (HLS) toolchain and integrated in the
Vivado Design Suite.

\textbf{Workload.} The evaluation uses the LFM2.5-1.2B language backbone (input
dimension 2048, feed-forward dimension 8192), which is shared across the Instruct and
Thinking fine-tuned variants and the base version. The accelerator processes one token
per invocation, with multi-token prefill handled by sequential IP calls in the driver
loop, issued by the modified \texttt{llama.cpp}. This reflects the single-token
synthesis constraint imposed by on-chip memory limits.

\textbf{Baseline.} The baseline corresponds to CPU Only execution of the same model
through an identical inference pipeline. In the accelerated configuration, only the
MatMul + SwiGLU block is offloaded to the FPGA, while all remaining operations are
executed on the CPU. Both configurations are produced by the same modified
\texttt{llama.cpp} binary, with the offload enabled or disabled at runtime through the
\texttt{LLAMA\_SWIHW} environment variable (\texttt{1} for FPGA+CPU, \texttt{0} for
CPU Only), without recompilation. The CPU Only baseline uses the standard Q4\_K\_M
distribution, whereas the FPGA+CPU path uses the uniformly re-quantized Q4\_K model that
the accelerator requires (Table~\ref{tab:short_results}).

\section{Evaluation Metrics}
\label{sec:metrics}

The evaluation captures both performance and efficiency through a set of directly
measured and derived quantities. Throughput, expressed in tokens per second, is measured
separately for the prefill and decode phases using \texttt{llama-bench} with a fixed
workload of 12 prompt tokens and 64 generated tokens, repeated over 10 runs to obtain
stable numbers. Power consumption is recorded concurrently via the onboard telemetry
interface (\texttt{xmutil platformstats}, an integrated system monitoring utility on the
KV260 platform), which samples the system-on-module total power at 1\,Hz throughout each
benchmark session. From these two measurements, performance per watt (tokens/s/W) is
derived as the ratio of throughput to mean power, providing a joint view of speed and
energy efficiency across operating points. FPGA resource utilization, reported in terms
of look-up tables (LUT), flip-flops (FF), block RAM (BRAM), ultra RAM (URAM), and DSP
slices, is taken from the post-place-and-route reports produced by the Vivado
implementation flow and quantifies how closely the accelerator design approaches the
physical limits of the target device. Together, these metrics characterize the
trade-off space between throughput, efficiency, and resource utilization that governs
the feasibility of edge deployment.

\section{Results}
\label{sec:results}

This section compares CPU Only and FPGA+CPU execution across thread counts (T1--T4)
using throughput and performance-per-watt. Table~\ref{tab:short_results} reports the
main quantitative results, while Fig.~\ref{fig:utilization} provides hardware resource
usage.

\begin{table}[htbp]
\centering
\scriptsize
\caption{Throughput and efficiency across platforms and thread counts. Bold marks the
preferable value per metric within each KV260 thread-count pair.}
\label{tab:short_results}
\setlength{\tabcolsep}{3pt}
\begin{tabular}{c c c c c c c}
\hline
Threads & Mode & Prefill (t/s) & Decode (t/s) & Power (W) & Prefill (t/s/W) & Decode (t/s/W) \\
\hline
\multicolumn{7}{c}{\itshape {KV260 (4$\times$ Cortex-A53, CPU Only: Q4\_K\_M \textbar{} FPGA+CPU: Q4\_K, 250\,MHz FPGA fabric)}} \\
\hline
T1 & CPU Only & 1.09 & 0.88 & 3.96 & 0.275 & 0.222 \\
T1 & FPGA+CPU & \textbf{2.20} & \textbf{1.63} & 5.10 & \textbf{0.431} & \textbf{0.320} \\
\hline
T2 & CPU Only & 2.17 & 1.66 & 4.25 & 0.511 & 0.391 \\
T2 & FPGA+CPU & \textbf{2.90} & \textbf{2.21} & 5.15 & \textbf{0.563} & \textbf{0.429} \\
\hline
T3 & CPU Only & \textbf{3.16} & \textbf{2.36} & \textbf{4.41} & \textbf{0.717} & \textbf{0.535} \\
T3 & FPGA+CPU & 3.06 & 2.05 & 5.14 & 0.595 & 0.399 \\
\hline
T4 & CPU Only & \textbf{3.52} & \textbf{2.31} & \textbf{4.28} & \textbf{0.822} & \textbf{0.540} \\
T4 & FPGA+CPU & 2.85 & 1.46 & 5.13 & 0.556 & 0.285 \\
\hline
\multicolumn{7}{c}{\textit{STM32MP257F-DK (2$\times$ Cortex-A35, Q4\_K\_M, power analytically estimated\textsuperscript{b})}} \\
\hline
T1 & CPU Only & 0.83 & 0.67 & 0.62 & 1.339 & 1.081 \\
T2 & CPU Only & 1.61 & 1.28 & 0.80 & 2.013 & 1.600 \\
\hline
\end{tabular}

\par\vspace{4pt}
{\footnotesize
\textsuperscript{b}Derived from STM32MP257 datasheet characterization~\cite{ds14284}
(Tables~22--26) and PMIC rail voltages~\cite{um3385} (\S7.5.3); see
Section~\ref{subsec:stm32comparison} for full derivation.}
\end{table}

Peak resident memory is approximately 728\,MB for CPU execution (the GGUF model is
memory-mapped and demand-paged; measured peak Resident Set Size 728\,MB across all
thread counts) and approximately 1.4\,GB for the FPGA+CPU path (the same 728\,MB process
RSS plus a 640\,MB CMA-allocated udmabuf DMA buffer that does not appear in RSS).
Per-token decode latency (the reciprocal of decode throughput) ranges from 1136\,ms at
T1 CPU Only down to 424\,ms at T3 and 433\,ms at T4; for FPGA+CPU it is 613\,ms at T1,
452\,ms at T2, 488\,ms at T3, and 685\,ms at T4.

\begin{figure}[!ht]
  \centering
  \includegraphics[width=0.65\linewidth]{utilization_final.png}
  \caption{Utilization of the Kria KV260 on-board resources for the implementation of the feed-forward block accelerator.}
  \label{fig:utilization}
\end{figure}

\section{Comparison with STM32MP257F-DK}
\label{subsec:stm32comparison}

To contextualize the KV260 results against a different Arm-class edge platform, the
standard LFM2.5-1.2B Q4\_K\_M distribution (8 of 16 down-projection layers at Q6\_K, the
remainder at Q4\_K) was benchmarked on the STM32MP257F-DK, which integrates on chip two
Cortex-A35 cores running at up to 1.5\,GHz. Results are reported alongside the KV260 data
in Table~\ref{tab:short_results}. Total board power consumption is derived analytically
as the sum of four terms:
\begin{enumerate}[label=(\roman*)]
    \item \textbf{SoC rail power}: datasheet characterization currents~\cite{ds14284} (Tables~22--26, $T_j = 25\,^\circ\text{C}$) multiplied by the PMIC output voltages programmed in the board user manual~\cite{um3385} (\S7.5.3, $V_{\text{DDCPU}} = 0.91\,\text{V}$, $V_{\text{DDCORE}} = 0.82\,\text{V}$); this yields $174\,\text{mW}$ at T1 and $320\,\text{mW}$ at T2,
    \item \textbf{LPDDR4 power}: background-dominated at fewer than 5\% utilization ($\approx 150\,\text{mW}$ independent of thread count),
    \item \textbf{Board peripheral idle power} (GigE PHY, STLINK, Wi-Fi module, eMMC): $\approx 205\,\text{mW}$, constant across T1 and T2,
    \item \textbf{PMIC buck conversion losses} at the operating load, assuming $\eta = 0.86$ for the SoC/DRAM rails and $\eta = 0.88$ for the $3.3\,\text{V}$ peripheral rail.
\end{enumerate}
Summing all terms gives a total board input power of $0.62\,\text{W}$ (T1) and
$0.80\,\text{W}$ (T2).

\section{Discussion}
\label{sec:discussion}

The FPGA+CPU solution improves efficiency at lower thread counts but not uniformly. At
T1, performance-per-watt improves by +57\% (prefill) and +44\% (decode) versus CPU Only,
with board power rising from 3.96\,W (CPU) to 5.10\,W (FPGA+CPU). At T2, the gains narrow
to +10\% for both phases (4.25\,W versus 5.15\,W). At T3 and T4, CPU Only is superior:
FPGA+CPU is lower by 17\%/25\% at T3 and 32\%/47\% at T4 (prefill/decode). At T4,
FPGA+CPU decode degrades further to 1.46\,t/s (below the T3 value of 2.05\,t/s), because
all four A53 threads compete with the UIO interrupt handler for scheduler time,
increasing per-token round-trip latency. These trends are consistent with end-to-end
throughput being constrained more by orchestration and memory movement than by
arithmetic scaling alone.

The results indicate that offloading the MatMul + SwiGLU feed-forward block to the FPGA
is beneficial because it targets a runtime-dominant kernel, consistently with prior
transformer accelerators that focus on FFN/MLP and attention substructures as the most
advantageous hardware targets~\cite{lu2020hw-mha}. The accelerator is effective when
host-side contention is limited and the cost of orchestration can be amortized
(T1--T2); at higher thread counts (T3--T4), the incremental benefit is offset by
transfer overhead and synchronization, and the CPU alone can supply and process the
residual stream more efficiently.

Comparing the KV260 against the STM32MP257F-DK isolates the effect of the wider
Cortex-A53 pipeline and FPGA fabric from the Cortex-A35 baseline. The A35 achieves
1.28\,t/s decode at T2, versus 1.66\,t/s for the KV260 A53 at T2 (CPU Only) and
\textbf{2.21}\,t/s for the KV260 FPGA+CPU path. This places the KV260 FPGA+CPU
configuration at a \textbf{1.7$\times$} decode speedup over the A35 at matched thread
count, at the cost of a more complex integration and approximately \textbf{4.4}\,W of
additional board power (5.15\,W FPGA+CPU versus 0.80\,W for the STM32 at T2). The
STM32's 3.7$\times$ efficiency advantage (1.60\,t/s/W decode versus 0.429\,t/s/W for the
KV260 FPGA+CPU) confirms that the accelerator is justified by its latency margin: the
1.7$\times$ decode speedup brings the system within the 452\,ms real-time window that
the A35 at 781\,ms per token cannot reach.

Section~\ref{sec:research_question} defined a per-token generation latency target of
200--500\,ms for the language model component. Of the eight KV260 configurations, four
enter this range: T3 and T4 CPU Only (424 and 433\,ms) and T2 and T3 FPGA+CPU (452 and
488\,ms). No configuration reaches the 200\,ms instantaneous threshold. The FPGA+CPU path
achieves sub-500\,ms decode with only two active CPU threads, whereas CPU Only requires
three or four to enter the same latency range. This thread advantage matters in practice:
a full conversational pipeline must concurrently run ASR and TTS stages that compete for
the same A53 cores, so reducing the language model's thread demand directly widens the
scheduling budget for the surrounding stages. At the same time, the 452\,ms LM-only
latency at T2 FPGA+CPU is already a significant fraction of the 500\,ms budget, leaving
little headroom for ASR and TTS; meeting the 0.5\,s aggregate target in a streaming
pipeline will require either a faster platform or accepting the ``immediate'' perceptual
range (0.5--1\,s) defined in Section~\ref{sec:research_question}.

The contrast highlights that the practical value of the FPGA offload scales with the gap
between the host CPU and the accelerator: on a weaker host the absolute gain is larger,
while on a stronger host (more A53 threads) the orchestration overhead increasingly
dominates. Table~\ref{tab:kv260_comparison} places this design within the space of
published KV260 accelerators; unlike compute-intensive matrix-multiply designs that
saturate DSP slices, the MatMul + SwiGLU kernel saturates routing fabric (87\% LUT,
20\% DSP), reflecting a load-dominated rather than arithmetic-dominated workload.

\begin{table}[H]
\centering
\caption{KV260 accelerator design comparison.}
\label{tab:kv260_comparison}
\setlength{\tabcolsep}{2.5pt}
\footnotesize
\begin{tabular}{l c c c c}
\hline
 & \textbf{Tiny Transf.} & \textbf{Self-Attn.} & \textbf{MatMul Accel.} & \textbf{This Work} \\
 & \cite{11310944tin} & \cite{li2025fpgabasedhardware} & \cite{11311069} & \\
\hline
Accelerated function & Full E2E pipeline & Q/K/V projections & Matrix mult. & MatMul + SwiGLU \\
Model / matrix size & 6.6k params. & DistilBERT & DistilBERT & LFM2.5 1.2B \\
Model size (params)  & 6.6k & 66M & 66M & 1.2B \\
Precision & float32 & INT8 & INT8 & INT8 / INT4 \\
Clock (MHz) & 77 & 100 & 100 & 250 \\
\hline
Ops/cycle (INT8 MAC) & --- & 1024 & 1024 & 32 \\
Pure compute throughput (GOPS)\textsuperscript{d}      & --- & 3.12\textsuperscript{c} & 22.83\textsuperscript{c} & 8.00\textsuperscript{a} \\
Effective system throughput (GOPS)\textsuperscript{e} & --- & 2.85 & 13.52 & 2.58\textsuperscript{b} \\
\hline
LUT (\%)          & 30 & 60 & 84 & 87 \\
DSP (\%)          & 27 & 83 & 83 & 20 \\
BRAM (\%)         & 37 & 88 & 71 & 32 \\
FF (\%)           & 20 & 43 & 44 & 50 \\
URAM (\%)         & --- & --- & --- & 13 \\
\hline
Power (W)         & --- & 3.3 & --- & 4.8--6.0 \\
\hline
\end{tabular}

\par\vspace{4pt}
{\footnotesize
\textsuperscript{a}32 parallel INT32 MAC chains (8 blocks $\times$ 4 rows) at 250\,MHz.
\textsuperscript{b}Stage~5 (down projection).
\textsuperscript{c}Evaluation uses an FFN matrix multiplication case with dimensions $(64,768)\times(768,3072)$.
\textsuperscript{d}Hardware-only peak; excludes memory transfer and host CPU overhead.
\textsuperscript{e}Includes all system-level effects: memory transfer, host CPU communication, and orchestration.}
\end{table}

Results reported using larger FPGA platforms suggest that greater memory bandwidth and
fabric capacity help kernel-level gains and end-to-end speedup~\cite{FTRANS}. Direct
comparison with the 6.6k-parameter Tiny Transformer on the same platform highlights the
resource intensity of conversational architectures. While the Tiny Transformer maps its
entire end-to-end pipeline at low resource utilization across all fabric categories
(Table~\ref{tab:kv260_comparison}), accelerating just the feed-forward MatMul + SwiGLU
block of the 1.2B-parameter LFM2.5 model nearly saturates the same device at 87\% LUT and
32\% BRAM utilization at 250\,MHz. Critically, the Tiny Transformer's small parameter
count is not merely a resource advantage~\cite{11310944tin}: by fitting the entire
pipeline on chip, it eliminates off-chip weight fetching entirely, allowing execution to
remain compute-bound and sustain near-peak arithmetic throughput, a regime that larger
models cannot achieve on the same device. The Self-Attention
accelerator~\cite{li2025fpgabasedhardware} highlights differing architectural trade-offs:
its two-level tiling strategy keeps both matrix operands within on-chip BRAM across the
Q, K, and V projections, allowing a fully-unrolled $32\times32$ MAC array to sustain 1024
concurrent INT8 operations per cycle without off-chip bandwidth pressure, reaching 91\%
effective utilization of peak throughput. This is achievable because the DistilBERT-scale
projection matrices fit within the on-chip tiling footprint. The LFM2.5 SwiGLU weights,
by contrast, total approximately 30\,MB per layer across the three Q4\_K projections,
far beyond the ZU5EV's combined on-chip capacity of roughly 3.4\,MB (block, ultra, and
distributed RAM combined), which makes DDR streaming unavoidable and places the design in a
memory-bandwidth-limited regime regardless of MAC array width or tiling strategy. A
comparable BRAM-banked matrix multiplication study~\cite{11311069} reports similar severe
PS-side bottlenecks, reaching only 59\% of its peak throughput; the present work reaches
34--41\% of peak capacity, reflecting its higher per-token weight volume. The contrast in
DSP usage (83\% versus 20\%) reflects an operand-width mismatch rather than a deliberate
design choice: the MatMul + SwiGLU kernel's core MAC operation multiplies a 4-bit nibble
weight by an INT8 activation, yielding a 12-bit product too narrow to use a DSP48E2 slice
efficiently. The DSP48E2 is optimized for $18\times27$-bit multiplications, so HLS maps
the narrow multiply-accumulate operations to LUT-based logic instead; the 20\% DSP usage
in this work arises entirely from the FP32 scale and minimum-factor accumulation paths,
not from the dominant weight--activation inner loop.

Extending hardware acceleration beyond the MatMul + SwiGLU block to the Final Linear
Projection and Gated Short Convolution yields diminishing returns, as they account for
only 10.96\% and 10.44\% of total execution time, respectively. Furthermore, simultaneous
acceleration is physically infeasible on the KV260, as the MatMul + SwiGLU kernel alone
leaves insufficient resources for additional hardware mapping. To characterize the
attainable design space, roughly a dozen distinct architectural variants were synthesized
and evaluated during development, systematically varying the nibble storage width, the
BRAM banking strategy, the AXI burst organization, and the use of function inlining and
stream decoupling. Table~\ref{tab:design_space} summarizes the most informative attempts.
Their consistency is itself the main finding: three independent ceilings emerged, and
every variant that attacked one of them converged on the same wall.

\begin{table}[H]
\centering
\caption{Representative architectural variants explored during development and the ceiling
each one revealed.}
\label{tab:design_space}
\resizebox{\textwidth}{!}{%
\begin{tabular}{lll}
\hline
\textbf{Strategy} & \textbf{Result} & \textbf{Ceiling revealed} \\ \hline
Wide-BRAM nibble storage (gate/up) & $II=1$ load, 67 cycles/row; enables $K=4$ within budget & final choice \\
URAM $\rightarrow$ BRAM swap & No change & Four-write limit is operation-count, not latency \\
Cyclic-4 banking (proves independence) & No change & Four-write limit, not bank aliasing \\
Stream-decoupled dataflow & No change & Four-write fanout, not the AXI engine \\
Merged 4-row burst (gate/up) & $-38\%$ cycles but $+14$K LUT $\rightarrow$ 138K (118\%) & AXI/LUT budget: place-and-route fails \\
Wide-BRAM output load & $-15\%$ cycles, 74K routing overlaps & ZU5EV routing capacity \\
Full output-MAC unroll & $+102\%$ LUT, negligible speedup & ZU5EV routing/LUT capacity \\
\hline
\end{tabular}}
\end{table}

The first ceiling is AXI burst serialization. Each DDR burst pays a fixed setup overhead
of approximately 45 cycles on the ZU5EV's HP-port interconnect before data begins to flow,
so a single contiguous burst at an initiation interval above one outperforms several
parallel bursts at $II=1$. A merged load that fetches all four rows of Stages~2a/2b in one
320-word burst does reduce the projection-stage cycle count by 38\%, but the four-way
header-routing logic it requires adds about 14\,K LUTs, pushing the design to 138\,K
(118\% of the device) and failing place-and-route. The single-burst organization is
therefore retained, and the down-projection, whose rows are four times longer, cannot
apply the merge at all.

The second ceiling is the HLS scheduler's refusal to issue four memory writes in one
pipeline stage. Because the down-projection stores its nibbles in narrow 32-bit tiles
rather than the wide-BRAM layout used elsewhere (Section~\ref{subsec:predecode}), each
128-bit DDR word must be split and written into four separate single-port buffers within
the same cycle. Vitis HLS will not schedule four writes per cycle, and, crucially,
this is a limit on the \emph{number} of memory operations, not a physical port conflict: a
cyclic partitioning that proves the four banks are mutually independent does not help, and
swapping the buffers between BRAM and URAM changes nothing. The gate and up projections
escape the limit by storing each pre-decoded 128-bit word verbatim, a single write per
cycle; the down-projection cannot, because its 32 blocks per row (the 8192-element
contraction dimension divided by the 256-weight Q4\_K block) demand 32 simultaneous block,
header, and scale reads in the MAC, a density the ZU5EV fabric cannot route in the
wide-BRAM form.

That routing density is the third ceiling. A fully-unrolled output MAC array, fed by
wide-BRAM reads, gate-cache lookups, and sub-block-scale accesses, exceeds the device's
routable interconnect at 250\,MHz: the wide-BRAM output variant reaches its target cycle
count in synthesis but leaves more than 74,000 unresolved routing-node overlaps in
implementation. A migration to the ZU7EV (the next device in the same Zynq UltraScale+
family, with approximately 230\,K LUTs on the same processor--fabric interface) would
provide the headroom to route the fully-unrolled array and to merge the load calls that
currently serialize the AXI path.

Taken together, these ceilings bound the achievable throughput at approximately 3.6\,t/s,
well within the Amdahl ceiling of $\sim$8\,t/s established in
Section~\ref{subsec:cpuexec}; the deployed design attains 61\% of this 3.6\,t/s design
ceiling. They reflect the structural mismatch between a memory-bandwidth-limited workload
and a device optimized for logic-dense, compute-bound kernels.

Table~\ref{tab:operating_points} consolidates the latency compliance, efficiency
outcome, and dominant limiting factor for each operating point discussed above, while
Table~\ref{tab:design_constraints} summarizes the architectural constraints that bound
the throughput ceiling independent of thread count.

\begin{table}[H]
\centering
\scriptsize
\caption{Operating-point summary: per-token decode latency, real-time budget
compliance, efficiency outcome, and dominant limiting factor per operating point.
Efficiency for the KV260 rows is decode t/s/W relative to CPU Only at the same thread
count; for the STM32 it is relative to KV260 FPGA+CPU at
T2\protect\textsuperscript{a}.}
\label{tab:operating_points}
\setlength{\tabcolsep}{3pt}
\begin{tabular}{c c r c r p{2.8cm}}
\hline
Threads & Mode & Latency (ms) & $<$500\,ms & Efficiency & Limiting factor \\
\hline
\multicolumn{6}{c}{\itshape KV260 (4$\times$ Cortex-A53)} \\
\hline
T1 & CPU Only  & 1136 & No           & baseline  & --- \\
T1 & FPGA+CPU  & 613  & No           & +44\%     & Single-thread CPU throughput ceiling \\
T2 & CPU Only  & 602  & No           & baseline  & --- \\
T2 & FPGA+CPU  & \textbf{452}  & \textbf{Yes} & +10\%  & AXI burst serialization \\
T3 & CPU Only  & \textbf{424}  & \textbf{Yes} & baseline  & --- \\
T3 & FPGA+CPU  & \textbf{488}  & \textbf{Yes} & $-$25\%   & CPU parallelism offsets FPGA \\
T4 & CPU Only  & \textbf{433}  & \textbf{Yes} & baseline  & --- \\
T4 & FPGA+CPU  & 685  & No           & $-$47\%   & UIO interrupt contention \\
\hline
\multicolumn{6}{c}{\itshape STM32MP257F-DK (2$\times$ Cortex-A35)} \\
\hline
T2 & CPU Only  & 781  & No           & $3.7{\times}$ more efficient\textsuperscript{a} & A35 core throughput \\
\hline
\end{tabular}
\par\vspace{4pt}
{\footnotesize\textsuperscript{a}Decode t/s/W relative to KV260 FPGA+CPU T2: 1.60 vs.\ 0.429\,t/s/W.}
\end{table}

\begin{table}[H]
\centering
\footnotesize
\caption{Design-level architectural constraints that bound the throughput ceiling
independent of thread count.}
\label{tab:design_constraints}
\setlength{\tabcolsep}{5pt}
\begin{tabular}{p{8.3cm} c p{5.4cm}}
\hline
\textbf{Constraint} & \textbf{Stage} & \textbf{Ceiling impact} \\
\hline
Per-burst AXI setup overhead ($\approx$45 cycles) makes many short bursts slower than
fewer large ones; a merged burst mitigates it in Stages~2a/2b but is infeasible in
Stage~5 & 2a, 2b, 5 & DDR utilization capped at 38--41\% of peak \\
\hline
The memory scheduler cannot write to four nibble-storage banks within the same pipeline
clock cycle, regardless of storage technology or partitioning strategy & 5 & Stage~5 is
load-bound; down-projection floored at 1.62M cycles \\
\hline
The ZU5EV routing fabric cannot simultaneously carry the 128-bit weight reads,
gate-cache lookups, and scale-factor accesses of a fully-unrolled output MAC array at
250\,MHz & 5 & Prevents the fully-unrolled output MAC array \\
\hline
\end{tabular}
\end{table}
